\subsection{Seudocódigo}

El seudocódigo está basado en lenguajes de programación, pero omite detalles de sintaxis (e.g., corchetes, dos puntos) o de implementación (e.g., modularidad, manejo de errores) de los lenguajes de programación que no afectan el algoritmo.

No existe una forma estándar de escribir seudocódigo: distintas fuentes utilizan convenciones diferentes inspiradas en distintos lenguajes de programación.

En estos apuntes se utiliza un seudocódigo inspirado en Python, por ser un lenguaje con una sintaxis breve y legible. Como la idea es describir el algoritmo, no se usan modismos específicos de Python: se escribe código familiar para cualquiera que domine algún lenguaje de programación imperativo.

\subsubsection{Ejemplo de seudocódigo}

La siguiente función recibe un arreglo de números y retorna el mayor:
\begin{pseudocode}
def find_largest(nums: arreglo[int]) -> int:
    largest = |\(-\infty\)|
    for i=0 to len(nums)-1:
        if nums[i] > largest:
            largest = nums[i]
    return largest
\end{pseudocode}

Es evidente que el código se asemeja a Python. Los detalles más recurrentes en estas notas donde el seudocódgo difiere con Python son los siguientes:
\begin{itemize}
  \item Al utilizar cualquier estructura de datos abstracta, como arreglos, colas, pilas, árboles, etc, la estructura se declara en español.
  \begin{itemize}
    \item E.g., se utiliza el tipo \inline{arreglo} en lugar de objetos de la clase \inline{list}. Eso pretende representar un arreglo abstracto, con tamaño fijo y acceso por índice en tiempo constante, en lugar del tipo \inline{list} de Python que está implementado como un arreglo dinámico.
  \end{itemize}
  \item Los ciclos sobre rangos de enteros se escriben como \inline{for i=0 to n-1} en lugar de \inline{for i in range(n)}, pues es más familiar para quienes provienen de otros lenguajes imperativos y además hace explícitos los límites del recorrido.
  \item En lugar de usar \textit{slicing}, se escribe literalmente (en español) que se realiza una copia y de qué segmento específico del arreglo o lista.
  \item Se prefiere la notación matemática sobre sintaxis específica de Python, por ejemplo se usa la lemniscata \(\infty\) en lugar de la sintaxis \verb|float('inf')|.
  \item Se evita usar cualquier función, como \inline{sorted}, o librería, como \inline{itertools}, que encapsule operaciones complejas, pues el propósito es precisamente comprender cómo se implementarían desde cero.
\end{itemize}

Así pues, debería ser trivial traducir cualquiera de los algoritmos en estas notas a Python y similarmente debería ser sencillo traducirlo a cualquier lenguaje imperativo.
